% Approved English source: EWAS_GSEG_TRS_Chapters_01-04_EN.md
% Ground Segment TRS; integrated product, WBS 123456.02.
\section{Scope}

Document \texttt{EWAS-HSB-GSEG-TRS-02} is the segment-level Technical Requirements Specification for the \textbf{EWAS Ground Segment (GSEG)}, WBS \texttt{123456.02}, designated \texttt{CI-02}. It defines functional, performance, interface, operational, and security requirements for the \textbf{integrated Ground Segment}. Detailed internal requirements for individual subsystems are addressed in lower-level specifications.

\subsection{Product Structure and Functional Allocation}

{\footnotesize\setlength{\tabcolsep}{4pt}\renewcommand{\arraystretch}{1.18}%
\arrayrulecolor{TableBorderBlue}\rowcolors{2}{white}{TableRowLightBlue}%
\begin{longtable}{|>{\raggedright\arraybackslash}p{\dimexpr0.23\textwidth-2\tabcolsep-1.33333\arrayrulewidth\relax}|>{\raggedright\arraybackslash}p{\dimexpr0.31\textwidth-2\tabcolsep-1.33333\arrayrulewidth\relax}|>{\raggedright\arraybackslash}p{\dimexpr0.46\textwidth-2\tabcolsep-1.33333\arrayrulewidth\relax}|}
\caption{EWAS Ground Segment Product Structure and Functional Allocation}\label{tab:gseg_wbs_scope}\\
\hline\rowcolor{TableHeaderBlue}
\thead{WBS} & \thead{Element} & \thead{Functional scope} \\ \hline
\endfirsthead
\hline\rowcolor{TableHeaderBlue}
\thead{WBS} & \thead{Element} & \thead{Functional scope} \\ \hline
\endhead
\texttt{123456.02} & \textbf{Ground Systems / Ground Segment} & Integrated EWAS ground infrastructure and functions. \\ \hline
\texttt{123456.02.01} & \textbf{Data Archiving System (DAS)} & Storage, protection, and controlled retrieval of mission data, analysis products, warnings, and operational records. \\ \hline
\texttt{123456.02.01.01} & Raw Data Archive & Raw data and received source products. \\ \hline
\texttt{123456.02.01.02} & Processed Data Archive & Processing results and derived data. \\ \hline
\texttt{123456.02.02} & \textbf{Command \& Control System (C\&C)} & Autonomous two-LEO selection using current orbital products, command preparation, validation, transmission coordination and execution monitoring. \\ \hline
\texttt{123456.02.03} & \textbf{Mission Control Center (MCC)} & Operator supervision, flight dynamics and mission planning, orbital/conjunction assessment, and post-warning operational actions. \\ \hline
\texttt{123456.02.04} & \textbf{Data Processing System (DPS)} & Observation correlation, data fusion, trajectory estimation, and impact-area assessment. \\ \hline
\texttt{123456.02.04.01} & Data Processing Software & Preparation and processing of input data. \\ \hline
\texttt{123456.02.04.02} & Event Detection Module & Ground-based event identification and assessment. \\ \hline
\texttt{123456.02.04.03} & Trajectory Estimation Module & Trajectory and impact-area estimation. \\ \hline
\texttt{123456.02.05} & \textbf{Ground Station Network (GSN)} & Satellite communications, telemetry and mission-data reception, and command transmission. \\ \hline
\texttt{123456.02.06} & \textbf{Early Warning System (EWS)} & Warning generation and delivery, including message status and subsequent versions. \\ \hline
\texttt{123456.02.07} & \textbf{Ground Infrastructure Subsystem (GIS)} & Shared computing and internal networks, common time services, electrical power, cooling and infrastructure monitoring. \\ \hline
\end{longtable}}

GSEG comprises seven subsystems. The seventh, GIS, is a supporting subsystem rather than a mission-processing step. The GSEG DDF provides the product tree; \texttt{123456.02.07} is the proposed new WBS identifier, pending alignment with the controlled WBS baseline. Subsystem specifications define detailed technical requirements.

\subsection{Mission Scope and Segment Boundaries}

EWAS provides \textbf{global launch detection and tracking}, with priority given to the assessment of threats to the Baltic region (\texttt{R-049}). The Ground Segment is planned as an operational installation at \textbf{one geographical location}; the final site and equipment configuration are defined in the DDF.

\begin{sloppypar}
GSEG encompasses terrestrial communication, computing, and operational infrastructure, including common services provided by GIS. Spacecraft, constellation, and onboard ISL design belong to the peer Space Segment (\texttt{123456.03}). EWAS responsibility for warning delivery terminates at the authorized authority interface; subsequent public alerting is outside the system boundary (\texttt{R-048}). The system does not provide weapon guidance or interceptor-control functions (\texttt{R-001}, \texttt{R-002}).
\end{sloppypar}

\section{Introduction}

The Early Warning Alarm System (EWAS) is an early warning system for ballistic missile launches, providing threat-assessment information to authorized crisis-management and civil-protection authorities. The system consists of the Space Segment (SSEG), which detects and tracks events, and the Ground Segment (GSEG), which supports mission operations, data analysis, and warning dissemination. This specification defines the GSEG requirements derived from system requirements \texttt{R-001}--\texttt{R-049}.

\subsection{Position in the Documentation Hierarchy}

The GSEG and SSEG specifications are peer documents that derive from system requirements and coordinate shared interfaces. GSEG requirements are further allocated to subsystems and their technical specifications.

Product documentation also includes the \textbf{DDF} (Design Definition File), which defines the design, and the \textbf{DDJF} (Design Justification File), which records the supporting analyses and design rationale.

\subsection{Specification Structure}

{\footnotesize\setlength{\tabcolsep}{4pt}\renewcommand{\arraystretch}{1.18}%
\arrayrulecolor{TableBorderBlue}\rowcolors{2}{white}{TableRowLightBlue}%
\begin{longtable}{|>{\raggedright\arraybackslash}p{\dimexpr0.39\textwidth-2\tabcolsep-1.5\arrayrulewidth\relax}|>{\raggedright\arraybackslash}p{\dimexpr0.61\textwidth-2\tabcolsep-1.5\arrayrulewidth\relax}|}
\caption{Ground Segment TRS Structure}\label{tab:gseg_document_structure}\\
\hline\rowcolor{TableHeaderBlue}
\thead{Part} & \thead{Content} \\ \hline
\endfirsthead
\hline\rowcolor{TableHeaderBlue}
\thead{Part} & \thead{Content} \\ \hline
\endhead
References & Applicable and reference documents, definitions, and abbreviations. \\ \hline
Scope and Introduction & Product scope, WBS, boundaries, and specification conventions. \\ \hline
3. User's Need Presentation & User needs and operational constraints. \\ \hline
4. Selected Concept / Product Presentation & Logical architecture and interaction of segment functions. \\ \hline
5. Life Profile Description & GSEG life cycle. \\ \hline
6. Environment and Constraints Description & Operating environment and constraints. \\ \hline
7. Technical Requirements & Unambiguous, verifiable GSEG requirements. \\ \hline
Annex A --- RTM & Traceability between system and segment requirements and their WBS allocations. \\ \hline
\end{longtable}}

The specification follows the structure of \texttt{ECSS-E-ST-10-06C} Annex A. Verification planning and recording of results are governed by the VP (\texttt{AD-11}) and VCD (\texttt{AD-12}).

\subsection{Requirement Conventions}

In Section 7, \textbf{shall} denotes a mandatory, verifiable requirement. \textbf{Should} denotes a recommendation, \textbf{may} denotes permission, while \textbf{can} and the descriptive present tense convey information rather than new requirements.

Requirement identifiers follow \texttt{GSEG-REQ-<CAT>-<NNN>}, where \texttt{CAT} is the category and \texttt{NNN} is a persistent three-digit number.

{\footnotesize\setlength{\tabcolsep}{4pt}\renewcommand{\arraystretch}{1.18}%
\arrayrulecolor{TableBorderBlue}\rowcolors{2}{white}{TableRowLightBlue}%
\begin{longtable}{|>{\raggedright\arraybackslash}p{\dimexpr0.16\textwidth-2\tabcolsep-1.5\arrayrulewidth\relax}|>{\raggedright\arraybackslash}p{\dimexpr0.84\textwidth-2\tabcolsep-1.5\arrayrulewidth\relax}|}
\caption{Ground Segment Requirement Categories}\label{tab:gseg_req_categories}\\
\hline\rowcolor{TableHeaderBlue}
\thead{Category} & \thead{Scope} \\ \hline
\endfirsthead
\hline\rowcolor{TableHeaderBlue}
\thead{Category} & \thead{Scope} \\ \hline
\endhead
FUN & Integrated segment functions \\ \hline
ARC & Architecture and configuration \\ \hline
PERF & Performance and response time \\ \hline
IF & Segment interfaces \\ \hline
DATA & Data processing, exchange, and retention \\ \hline
SEC & Information and access security \\ \hline
REL & Reliability and availability \\ \hline
OPS & Operations and maintenance \\ \hline
VER & Verification-related requirements \\ \hline
\end{longtable}}

Requirements are traced to their parent requirements in the RTM. Verification methods defined in \texttt{AD-04} are \textbf{Test (T), Analysis (A), Review of Design (R)}, and \textbf{Inspection (I)}.
